% GENERATED FILE. Edit canonical lessons, then run npm run generate:books.
% canonical-source-hash: fnv1a64:29e2569758b3fdc0
% canonical-lessons: JA-C441-jiman-suru, JA-C441-namaiki, JA-C441-erai, JA-C441-namakemono, JA-C441-fuzakeru

\chapter{Proud, Cheeky or Lazy}
\label{ch:ja-proud-cheeky-or-lazy}
\hlchaptermodality{441}

\begin{chapteropening}
\textbf{By the end of this chapter:} \emph{I can describe someone who boasts, slacks or fools around.}

Complete the last lesson of chapter 441: I can describe someone who boasts, slacks or fools around.
\end{chapteropening}

\section[\texorpdfstring{\ja{じまんする} (jiman suru)}{じまんする (jiman suru)}]{\ja{じまんする} (jiman suru) — to boast; to be proud of}

\label{lesson:JA-C441-jiman-suru}



\begin{quote}
Before the new one: say the Japanese for neatly, then the Japanese for suitable.
\end{quote}

\subsection*{You'll want to know: \ja{じまんする}}
\textbf{\ja{じまんする}} — \emph{jiman suru} — \textquotedblleft{}to boast; to be proud of\textquotedblright{}.

\begin{quote}
\textbf{\ja{ちちは} \ja{むすこを} \ja{じまんしています。}} — \emph{chichi wa musuko o jiman shite imasu} — \textquotedblleft{}My father is proud of his son.\textquotedblright{}
\end{quote}

Polite \textbf{\ja{じまんします}}, \emph{jimanshimasu}; a \ja{する} verb. Past (polite): \textbf{\ja{かれは} \ja{くるまを} \ja{じまんしました。}}, \emph{kare wa kuruma o jiman shimashita}.

\begin{cousinweb}
\textbf{\ja{じまんする}} (\emph{jiman suru}) is Sino-Japanese (kango) \emph{jiman} with the native verb \emph{suru}, \textquotedblleft{}to do\textquotedblright{}: two Chinese characters read with their on'yomi, readings borrowed from Chinese. They are \emph{ji}, \textquotedblleft{}oneself\textquotedblright{} (Middle Chinese \emph{dzijH}) and \emph{man}, \textquotedblleft{}ridicule\textquotedblright{} (\emph{maenH}) (KANJIDIC2; Middle Chinese in Baxter's transcription).
\end{cousinweb}

\begin{grammarlens}[title={What you've built}]
One more word for describing what you lived and felt.
\end{grammarlens}

\subsection*{Your turn}
\begin{itemize}
\raggedright
  \item \emph{Say it:} \emph{jiman suru}
  \item \emph{Say it:} \emph{jiman suru}, once more
  \item \emph{Say it:} \emph{chichi wa musuko o jiman shite imasu} — My father is proud of his son.
  \item \emph{Recall:} say the Japanese for neatly, then the Japanese for suitable, then say \emph{jiman suru} again
\end{itemize}

\begin{tcolorbox}[breakable,colback=teal!4,colframe=teal!35!black,title={Before you move on}]
What does \ja{じまんする} mean? (\textquotedblleft{}To boast; to be proud of\textquotedblright{}.) Say it once more.
\end{tcolorbox}
\section[\texorpdfstring{\ja{なまいき} (namaiki)}{なまいき (namaiki)}]{\ja{なまいき} (namaiki) — cheeky; impertinent}

\label{lesson:JA-C441-namaiki}



\begin{quote}
Before the new one: say the Japanese for suitable, then the Japanese for to boast.
\end{quote}

\subsection*{You'll want to know: \ja{なまいき}}
\textbf{\ja{なまいき}} — \emph{namaiki} — \textquotedblleft{}cheeky; impertinent\textquotedblright{}.

\begin{quote}
\textbf{\ja{おとうとは} \ja{すこし} \ja{なまいきです。}} — \emph{otōto wa sukoshi namaiki desu} — \textquotedblleft{}My younger brother is a little cheeky.\textquotedblright{}
\end{quote}

A na-adjective: \ja{な} before a noun.

\begin{cousinweb}
\textbf{\ja{なまいき}} (\emph{namaiki}) is a mixed word: native \emph{nama}, \textquotedblleft{}raw\textquotedblright{}, joined to Sino-Japanese \emph{iki}, \textquotedblleft{}spirit\textquotedblright{}, from \emph{i}, \textquotedblleft{}mind\textquotedblright{} (Middle Chinese \emph{'iH}), and \emph{ki}, \textquotedblleft{}spirit\textquotedblright{} (\emph{khj+jH}), borrowed from Chinese (JMdict, KANJIDIC2; Middle Chinese in Baxter's transcription).
\end{cousinweb}

\begin{grammarlens}[title={What you've built}]
One more word for describing what you lived and felt.
\end{grammarlens}

\subsection*{Your turn}
\begin{itemize}
\raggedright
  \item \emph{Say it:} \emph{namaiki}
  \item \emph{Say it:} \emph{namaiki}, once more
  \item \emph{Say it:} \emph{otōto wa sukoshi namaiki desu} — My younger brother is a little cheeky.
  \item \emph{Recall:} say the Japanese for suitable, then the Japanese for to boast, then say \emph{namaiki} again
\end{itemize}

\begin{tcolorbox}[breakable,colback=teal!4,colframe=teal!35!black,title={Before you move on}]
What does \ja{なまいき} mean? (\textquotedblleft{}Cheeky; impertinent\textquotedblright{}.) Say it once more.
\end{tcolorbox}
\section[\texorpdfstring{\ja{えらい} (erai)}{えらい (erai)}]{\ja{えらい} (erai) — admirable; great}

\label{lesson:JA-C441-erai}



\begin{quote}
Before the new one: say the Japanese for to boast, then the Japanese for cheeky.
\end{quote}

\subsection*{You'll want to know: \ja{えらい}}
\textbf{\ja{えらい}} — \emph{erai} — \textquotedblleft{}admirable; great\textquotedblright{}.

\begin{quote}
\textbf{\ja{まいにち} \ja{べんきょうして、えらいですね。}} — \emph{mainichi benkyō shite, erai desu ne} — \textquotedblleft{}You study every day: well done.\textquotedblright{}
\end{quote}

An i-adjective.

\begin{cousinweb}
\textbf{\ja{えらい}} (\emph{erai}) is a native Japanese word (wago), inherited rather than borrowed from Chinese; its kanji spelling was attached to it later (JMdict).
\end{cousinweb}

\begin{grammarlens}[title={What you've built}]
One more word for describing what you lived and felt.
\end{grammarlens}

\subsection*{Your turn}
\begin{itemize}
\raggedright
  \item \emph{Say it:} \emph{erai}
  \item \emph{Say it:} \emph{erai}, once more
  \item \emph{Say it:} \emph{mainichi benkyō shite, erai desu ne} — You study every day: well done.
  \item \emph{Recall:} say the Japanese for to boast, then the Japanese for cheeky, then say \emph{erai} again
\end{itemize}

\begin{tcolorbox}[breakable,colback=teal!4,colframe=teal!35!black,title={Before you move on}]
What does \ja{えらい} mean? (\textquotedblleft{}Admirable; great\textquotedblright{}.) Say it once more.
\end{tcolorbox}
\section[\texorpdfstring{\ja{なまけもの} (namakemono)}{なまけもの (namakemono)}]{\ja{なまけもの} (namakemono) — a lazy person}

\label{lesson:JA-C441-namakemono}



\begin{quote}
Before the new one: say the Japanese for cheeky, then the Japanese for admirable.
\end{quote}

\subsection*{You'll want to know: \ja{なまけもの}}
\textbf{\ja{なまけもの}} — \emph{namakemono} — \textquotedblleft{}a lazy person\textquotedblright{}.

\begin{quote}
\textbf{\ja{おとうとは} \ja{なまけものです。}} — \emph{otōto wa namakemono desu} — \textquotedblleft{}My younger brother is a lazybones.\textquotedblright{}
\end{quote}

\begin{cousinweb}
\textbf{\ja{なまけもの}} (\emph{namakemono}) is a native Japanese compound (wago), inherited rather than borrowed from Chinese: \emph{namake}, \textquotedblleft{}idling (namakeru)\textquotedblright{} and \emph{mono}, \textquotedblleft{}person, thing\textquotedblright{} (JMdict, KANJIDIC2).
\end{cousinweb}

\begin{grammarlens}[title={What you've built}]
One more word for describing what you lived and felt.
\end{grammarlens}

\subsection*{Your turn}
\begin{itemize}
\raggedright
  \item \emph{Say it:} \emph{namakemono}
  \item \emph{Say it:} \emph{namakemono}, once more
  \item \emph{Say it:} \emph{otōto wa namakemono desu} — My younger brother is a lazybones.
  \item \emph{Recall:} say the Japanese for cheeky, then the Japanese for admirable, then say \emph{namakemono} again
\end{itemize}

\begin{tcolorbox}[breakable,colback=teal!4,colframe=teal!35!black,title={Before you move on}]
What does \ja{なまけもの} mean? (\textquotedblleft{}A lazy person\textquotedblright{}.) Say it once more.
\end{tcolorbox}
\section[\texorpdfstring{\ja{ふざける} (fuzakeru)}{ふざける (fuzakeru)}]{\ja{ふざける} (fuzakeru) — to fool around; to joke}

\label{lesson:JA-C441-fuzakeru}



\begin{quote}
Before the new one: say the Japanese for admirable, then the Japanese for a lazy person.
\end{quote}

\subsection*{You'll want to know: \ja{ふざける}}
\textbf{\ja{ふざける}} — \emph{fuzakeru} — \textquotedblleft{}to fool around; to joke\textquotedblright{}.

\begin{quote}
\textbf{\ja{じゅぎょうで} \ja{ふざけないで} \ja{ください。}} — \emph{jugyō de fuzakenaide kudasai} — \textquotedblleft{}Please do not fool around in class.\textquotedblright{}
\end{quote}

Polite \textbf{\ja{ふざけます}}, \emph{fuzakemasu}; an ichidan verb. Past (polite): \textbf{\ja{がくせいが} \ja{ふざけました。}}, \emph{gakusei ga fuzakemashita}.

\begin{cousinweb}
\textbf{\ja{ふざける}} (\emph{fuzakeru}) is a native Japanese word (wago), inherited rather than borrowed from Chinese; JMdict writes it in kana (JMdict).
\end{cousinweb}

\begin{grammarlens}[title={What you've built}]
One more word for describing what you lived and felt.
\end{grammarlens}

\subsection*{Your turn}
\begin{itemize}
\raggedright
  \item \emph{Say it:} \emph{fuzakeru}
  \item \emph{Say it:} \emph{fuzakeru}, once more
  \item \emph{Say it:} \emph{jugyō de fuzakenaide kudasai} — Please do not fool around in class.
  \item \emph{Recall:} say the Japanese for admirable, then the Japanese for a lazy person, then say \emph{fuzakeru} again
\end{itemize}

\begin{tcolorbox}[breakable,colback=teal!4,colframe=teal!35!black,title={Before you move on}]
What does \ja{ふざける} mean? (\textquotedblleft{}To fool around; to joke\textquotedblright{}.) Say it once more.
\end{tcolorbox}
